%%%%%%%%%%  scan idx 219  |  folio 206  %%%%%%%%%%
quence du lemme suivant~:

\textbf{Lemme 2.2.1.} Soient $U$ un ouvert de $\R^m$, $V$ un ouvert convexe de $\R^n$ contenant $0$. Soient $A_0,A_1$ des fermés de $\{s\}\times U$, soit $U_0$ (resp. $U_1$) un ouvert de $S\times U$ (resp. $\{s\}\times U$) contenant $A_0$ (resp. $A_1$). Soit $f_1:U_1\to V$ un morphisme et soit $f_0:U_0\to S\times V$ un $S$-morphisme induisant sur $U_0\cap(\{s\}\times U)$ un morphisme définissant le même germe que $f_1$ au voisinage de $A_1$. Il existe alors un ouvert $U'$ de $S\times U$ contenant $s\times U$, et un $S$-morphisme $\bar f:U'\longrightarrow S\times V$ tel que
\begin{itemize}
\item[(i)] le germe de $\bar f$ suivant $\{s\}\times U$ prolonge le germe défini par $f_0$~;
\item[(ii)] le morphisme induit par $\bar f$ sur $\{s\}\times U$ définisse le même germe que $f_1$ au voisinage de $A_1$.
\end{itemize}

\emph{Preuve.} Quitte à remplacer $f_1$ par $af_1$, où $a$ est une fonction $C^r$ sur $U$, valant $1$ au voisinage de $A_1$ et à support dans $U_1$, on peut supposer que $U_1=U$. Soit $b$ une fonction $C^r$ sur $U$, valant $1$ sur $A_0$ et à support dans $U_0\cap(\{s\}\times U)$. Alors
\[
  U'=U_0\cup(S\times(U-\operatorname{Supp}b))
\]
est un ouvert de $S\times U$ contenant $\{s\}\times U$. Le morphisme $\bar f:U'\to S\times V$ défini par
\[
  \bar f(t,x)=\bigl(t,\ b(x)\pr_2f_0(t,x)+(1-b(x))f_1(x)\bigr)
\]
satisfait visiblement aux conditions (i) et (ii).

\textbf{Corollaire 2.3.} Sous les hypothèses de (2.2), les faisceaux $i_s^{-1}\underline{\Homop}(X,Y)$, $\underline{\Homop}(X_s,Y_s)$ sont mous. En particulier, les faisceaux $i_s^{-1}(\cO_X)$, $\cO_{X_s}$ sont mous\footnote{Pour $\cO_{X_s}$, c'est «~très bien connu~». Le rédacteur ignore si $\cO_X$ est mou quand les hypothèses de (2.2) sont vérifiées pour chaque $s\in S$ et que $X$ est paracompact. Cela impliquerait trivialement (moyennant lesdites hypothèses) la mollesse de $i_s^{-1}(\cO_X)$ ainsi que (2.5) ci-dessous.}. Pour tout fermé $A$ de $X_s$, l'application
\[
  H^0\bigl(X_s,i_s^{-1}\underline{\Homop}(X,Y)\bigr)
  \longrightarrow
  H^0\bigl(A,\underline{\Homop}(X_s,Y_s)\bigr)
\]
déduite de (2.1.1) est surjective.

